\section{USAMO 2007}
        \begin{enumerate}
        	\item (\textit{USAMO 2007}) Let $n$ be a positive integer. Define a sequence by setting $a_1=n$ and, for each $k>1$, letting $a_k$ be the unique integer in the range $0\le a_k\le k-1$ for which $a_1+a_2+\cdots+a_k$ is divisible by $k$. For instance, when $n=9$ the obtained sequence is $9, 1, 2, 0, 3, 3, 3, \ldots$.  Prove that for any $n$ the sequence $a_1, a_2, a_3, \ldots$ eventually becomes constant.


 

	\item (\textit{USAMO 2007}) A square grid on the Euclidean plane consists of all points $(m,n)$, where $m$ and $n$ are integers. Is it possible to cover all grid points by an infinite family of discs with non-overlapping interiors if each disc in the family has radius at least 5?


 

	\item (\textit{USAMO 2007}) Let $S$ be a set containing $n^2+n-1$ elements, for some positive integer $n$. Suppose that the $n$-element subsets of $S$ are partitioned into two classes. Prove that there are at least $n$ pairwise disjoint sets in the same class.


 

	\item (\textit{USAMO 2007}) An \textit{animal} with $n$ \textit{cells} is a connected figure consisting of $n$ equal-sized cells.${}^1$ The figure below shows an 8-cell animal.


 
                    \begin{figure}[H]
                        \centering
                        \includegraphics[width=0.5\textwidth]{Problem/USAMO_I_2007_4_Problem_Figure_1.PNG}
                    \end{figure}
                    A \textit{dinosaur} is an animal with at least 2007 cells. It is said to be \textit{primitive} if its cells cannot be partitioned into two or more dinosaurs. Find with proof the maximum number of cells in a primitive dinosaur.


 <small>${}^1$Animals are also called \textit{polyominoes}. They can be defined inductively. Two cells are \textit{adjacent} if they share a complete edge. A single cell is an animal, and given an animal with $n$ cells, one with $n+1$ cells is obtained by adjoining a new cell by making it adjacent to one or more existing cells.</small>


 

	\item (\textit{USAMO 2007}) Prove that for every nonnegative integer $n$, the number $7^{7^n}+1$ is the product of at least $2n+3$ (not necessarily distinct) primes.


 

	\item (\textit{USAMO 2007}) Let $ABC$ be an acute triangle with $\omega$, $\Omega$, and $R$ being its incircle, circumcircle, and circumradius, respectively. Circle $\omega_A$ is tangent internally to $\Omega$ at $A$ and tangent externally to $\omega$. Circle $\Omega_A$ is tangent internally to $\Omega$ at $A$ and tangent internally to $\omega$. Let $P_A$ and $Q_A$ denote the centers of $\omega_A$ and $\Omega_A$, respectively.  Define points $P_B$, $Q_B$, $P_C$, $Q_C$ analogously.  Prove that

 \[8P_AQ_A \cdot P_BQ_B \cdot P_CQ_C \le R^3,\]
with equality if and only if triangle $ABC$ is equilateral.


 

\end{enumerate}